\begin{enumerate}
	\item $\binom{14}{2}\times \binom{10}{2} = \np{4095} \neq 272$. 

\emph{L'énoncé indique \og{}~est-il possible~\fg{} et ne fait pas intervenir le nombre maximum de groupes composés de deux filles et deux garçons.
}
Comme $272 < \np{4095}$, \textbf{l'affirmation est vraie.}
	
	\item $f$ est dérivable comme composée de fonctions dérivables. 

$\forall x \in \R,~ f'(x) = 6\times \cos(2x+\pi)$	

$f'\left(\dfrac{\pi}{2}\right) = 6\times \cos\left(2\times\dfrac{\pi}{2} +\pi\right) = 6$

$f\left(\dfrac{\pi}{2}\right) = 3\times \sin\left(2\times\dfrac{\pi}{2} +\pi\right) = 0$
	
	L'équation de la tangente au point d'abscisse $\dfrac{\pi}{2}$ est donnée par :
	
	$y = \left(x - \dfrac{\pi}{2}\right)\times f'\left(\dfrac{\pi}{2}\right) + f\left(\dfrac{\pi}{2}\right)$
	
	$y = 6x- 6\times\dfrac{\pi}{2} $
	
	$y= 6x- 3\pi$
	
	\textbf{L'affirmation est vraie.}
	
	\item $\forall x\in ]0~;~+\infty[, F'(x) = \dfrac{2 x \ln x+2 x+1}{x}$, en particulier $F'(1) = \dfrac{3}{2}$ et $f(1) = 2$. $f$ n'est pas la dérivée de $F$. Ainsi on montre que $F$ n'est pas une primitive de $f$.
	
	\textbf{L'affirmation est fausse.}
	
	\item $g(0) =  45 \times \e^{0} +20 = 65$
	
	$\forall t \in \R\quad g'(t) +\np{0.06}\times g(t) = 45\times \np{0.06}\times \e^{ \np{0.06}t} + \np{0.06}\times\left(45\times\e^{ \np{0.06}t}+20\right) =  \np{5.4}\times\e^{ \np{0.06}t}+\np{1.2}$
	$g$ n'est pas solution de $(E)$.
	
	\textbf{L'affirmation est fausse.}
	
	\item Soit $y$ une solution positive de $\left(E_2\right)$ on peut écrire $\forall x \in \R\quad y'(x) = y(x) + 3\e^{\np{0.4}x}$, comme $\forall x \in \R\quad y(x)\geqslant 0\text{ et } 3\e^{\np{0.4}x}\geqslant 0$, alors  $\forall x \in \R\quad y'(x)\geqslant 0$.
	
	On dérive l'équation différentielle :
	
	$\forall x \in \R\quad y''(x) = y'(x) + \np{1.2}\e^{\np{0.4}x}$. Comme $\forall x \in \R\quad y'(x)\geqslant 0\text{ et } \np{1.2}\e^{\np{0.4}x}\geqslant 0$, alors  $\forall x \in \R\quad y''(x)\geqslant 0$.
	
	$\forall x \in \R\quad y''(x)\geqslant 0$ donc $y$ est convexe sur $\R$.
	
	\textbf{L'affirmation est vraie} : les solutions positives de $\left(E_2\right)$ sont des fonctions convexes sur $\R$.
\end{enumerate}
